\section{Related Work}
\subsection{Vision-Language-Action Models}

Vision--Language--Action (VLA) models extend pretrained vision--language models to robotic control by predicting actions conditioned on visual observations and language instructions~\citep{kim2025openvla, black2410pi0, black2025pi, bu2025univla, zhangdreamvla, cen2025worldvla, li2025controlvla, jia2026guidedvla}. OpenVLA-style methods~\citep{kim2025openvla, bu2025univla} formulate action generation as autoregressive prediction over discrete action tokens, whereas the $\pi$ series~\citep{black2410pi0, black2025pi} adopts continuous regression and flow matching for more precise and efficient action generation. Despite this progress, many VLA policies remain largely 2D-centric~\citep{zitkovich2023rt, huvideo, brohan2022rt}, limiting their ability to reason about the 3D structure required for complex manipulation. Recent methods~\citep{li2025bridgevla, fan2026any3d, qu2025spatialvla, sun2025geovla} incorporate explicit 3D information, such as depth or point clouds, to improve spatial understanding. However, these approaches often require additional 3D inputs at inference time and primarily focus on static geometry rather than 3D dynamics. ELAN4D addresses this limitation by introducing 4D-aware auxiliary supervision during training while keeping the policy interface unchanged at inference time.

\subsection{Predictive Supervision for Robotic Manipulation}

Recent studies improve robot policies through future prediction~\citep{hafnerdream, cen2025worldvla, xu2026futurevlajointvisuomotorprediction, jiang2025rynnvla001usinghumandemonstrations, fan2026futurevlaforecastingunifiedtrajectories, hafner2022deep, qian2026escapeepisodicspatialmemory}, including forecasting future RGB observations~\citep{wu2024unleashing, cen2025worldvla} or depth~\citep{zhangdreamvla}. However, these objectives are often 2D-centric and provide limited explicit spatial structure~\citep{cen2025worldvla}. Moreover, single-frame prediction may be insufficient to capture the physical dynamics needed for precise manipulation~\citep{du2023learning, black2023zero}. Recent methods such as Pri4R~\citep{kim2026pri4rlearningworlddynamics} and GeoPredict~\citep{qian2025geopredict} address this issue by using predictive supervision from 3D point tracks, i.e., 4D trajectories. Nevertheless, Pri4R relies on global tracks extracted by spatial trackers~\citep{xiao2025spatialtracker}, leading to a time-consuming preprocessing pipeline. GeoPredict injects predictive supervision into the VLM through additional track queries, coupling pretrained vision--language representations with low-level dynamics forecasting, which may interfere with the VLM's native capabilities~\citep{zhang2026vlmvla, driess2025knowledgeinsulatingvisionlanguageactionmodels}. In contrast, ELAN4D introduces embodiment-centric 4D supervision into the action expert through a ControlNet-style pathway with gradient isolation, preserving the original VLM representation while incurring negligible preprocessing cost.
\vspace{-1em}